\section{Positive contour measures and the short-range theorem}
\label{app:short-range-proof}

This appendix proves Theorem~\ref{thm:sr-bath}. Its source estimates,
phase definitions, and energy observable are kept explicit because a
signed partition-function approximation alone does not control a
positive reweighted probability measure.

The color threshold $q_0(d)$ is fixed once, here. It is the largest of:
the thresholds required by the cited finite-size scaling and contour
results \cite{BorgsKoteckyMiracleSole1991,BorgsChayesTetali2012} on
a fixed temperature interval around coexistence; the value $2d+1$,
used for the variance lower bound in
\eqref{eq:sr-variance-lower}; and the value needed for the
cluster-expansion criterion of \cite{BorgsChayesTetali2012} to hold
with the extra factor $e^{\rho m_\gamma}$ in
\eqref{eq:sr-cluster-majorant} for a fixed $\rho>0$. Every
``sufficiently large $q$'' below refers to this single threshold.

\subsection{Published finite-volume inputs and normalization}

We use the following established contour results, with source locations
given so that the deductions can be distinguished from their inputs.
Theorem 1 of Borgs, Koteck\'y, and Miracle-Sol\'e~\cite{BorgsKoteckyMiracleSole1991}
gives six-times differentiable dimensionless free energies
$\psi_o,\psi_d$ and a fixed real interval containing $\beta_c$ on which
\begin{equation}
 Z_N^{\mathrm{spin}}(\beta)
 =q e^{-N\psi_o(\beta)}+e^{-N\psi_d(\beta)}
   +O(e^{-bL})e^{-N\min_i\psi_i(\beta)}.
 \label{eq:sr-partition-input}
\end{equation}
The constants are uniform on a sufficiently small fixed compact interval.
The source states its error as $O(q^{-bL})$; for fixed $q$ this is
$O(e^{-bL})$ with a different constant $b$, which is the form used here.
At coexistence the free energies are equal and
\begin{equation}
 u_i=\psi_i'(\beta_c),\qquad
 u_o<u_d,\qquad \sigma_i^2=-\psi_i''(\beta_c).
 \label{eq:sr-phase-parameters}
\end{equation}
The ordered branch is stable above $\beta_c$, and the disordered branch
below. The source Hamiltonian in its equation (6), with $J=1$, is exactly
\eqref{eq:sr-hamiltonian}. The notation here uses $\psi_i=\beta f_i$
relative to that source. No complex-temperature continuation of
\eqref{eq:sr-partition-input} is assumed.

For the positive decomposition, put $p=1-e^{-\beta}$ and use bond sets
$A\subseteq\mathcal B_L$. Their random-cluster weights are
\begin{equation}
 w_\beta(A)=p^{|A|}(1-p)^{dN-|A|}q^{k(A)},
 \qquad
 Z_N^{\mathrm{RC}}(\beta)=\sum_Aw_\beta(A)
 =e^{-\beta dN}Z_N^{\mathrm{spin}}(\beta).
 \label{eq:sr-rc-normalization}
\end{equation}
This normalization follows by expanding
$\prod_{\{x,y\}}[1+(e^\beta-1)\ind_{\{\sigma_x=\sigma_y\}}]$.
It fixes the additive energy convention independently of conventions
used in contour references.

Throughout the short-range proof, contours, interiors and exteriors,
compatibility, phase events, restricted partition functions, and activities
are those of Borgs, Chayes, and Tetali~\cite{BorgsChayesTetali2012}
(BCT), Definitions 4.1--4.4 and 5.4, equations (6.1)--(6.5) and
(6.13)--(6.22). In particular, BCT defines the exterior of a contour
as the component containing a compatible interface, if such an interface
exists; otherwise it chooses the larger component, with a fixed tie rule.
This convention matters for large contours on the torus.

Geometrically, the construction thickens the occupied bonds and separates
ordered and disordered regions by their boundary components. Components
with zero winding vector are contours; the remaining components are
interfaces. In the absence of interfaces, the common exterior of the
contours has one label, ordered or disordered, which defines the positive
event $\mathcal O_L$ or $\mathcal D_L$. Configurations with at least one
interface form the tunneling event $\mathcal T_L$. Thus these events
partition the bond configurations and are independent of $\beta$.
Let $Z_{o,L}$ denote the ordered contribution with its color
multiplicity removed, and $Z_{d,L}$ the disordered contribution. BCT
(6.1)--(6.5) give exactly
\begin{equation}
 Z_N^{\mathrm{RC}}=q Z_{o,L}+Z_{d,L}+Z_{t,L}.
 \label{eq:sr-positive-contour-decomposition}
\end{equation}
At $\beta_c$, BCT Lemma 6.1 bounds the tunneling probability by
$C e^{-bL^{d-1}}$. The later construction
in~\cite{BorgsChayesHelmuthPerkinsTetali2022}, Definition 5,
uses a modified exterior convention; we do not identify its finite-volume
phase events with those used here or import estimates between them.

We will use the following specific parts of the original contour
construction. Write $m_\gamma=\|\gamma\|$ for the contour size.
Its diameter obeys
\begin{equation}
 \operatorname{diam}\gamma\le L,\qquad
 \operatorname{diam}\gamma\le m_\gamma/2,\qquad
 |\operatorname{Int}\gamma\cap\Lambda_L|
       \le\tfrac12m_\gamma\operatorname{diam}\gamma.
 \label{eq:sr-contour-geometry}
\end{equation}
These are the definition of diameter and Lemma 5.7
of~\cite{BorgsChayesTetali2012}; they hold for all $d\ge2$.
The positive restricted partition functions have the polymer form
\begin{equation}
 Z_{i,L}=e^{-e_i(\beta)N}
       \sum_{\Gamma\ \mathrm{compatible}}
       \prod_{\gamma\in\Gamma}K_i(\gamma;\beta).
 \label{eq:sr-polymer-form}
\end{equation}
Here the explicit ground coefficients are
$e_o=-d\log(1-e^{-\beta})$, $e_d=d\beta-\log q$, and
$\kappa=\tfrac12\log(e^\beta-1)$.
The activities are $e^{-\kappa(\beta)m_\gamma}$ times a ratio of
positive interior partition functions, with a factor $q$ in the
disordered case; see equations (6.18)--(6.22) of that source. The
coefficients $e_i$ and $\kappa$ have bounded derivatives on the fixed
temperature interval considered here.

Appendix A defines truncated activities
\begin{equation}
 K_i'(\gamma;\beta)=\min\{K_i(\gamma;\beta),
                                  e^{-\tau(\beta)m_\gamma}\},
 \label{eq:sr-truncation}
\end{equation}
where BCT (A.3) uses $\tau(\beta)=\beta/8-\beta/20+1$.
This positive exponential cap ensures uniform convergence of the cluster
expansion when $q$ is sufficiently large. Its infinite-volume
free energies will be denoted $f_i^{\mathrm{tr}}$, to distinguish them
from the $\psi_i$ in~\eqref{eq:sr-partition-input}. If
$a_i=f_i^{\mathrm{tr}}-\min_j f_j^{\mathrm{tr}}$, Lemma A.1(i) states
\begin{equation}
 a_i(\beta)\operatorname{diam}\gamma\le c_0\beta
 \quad\Longrightarrow\quad K_i(\gamma;\beta)=K_i'(\gamma;\beta).
 \label{eq:sr-inactive-cutoff-input}
\end{equation}
Moreover, equation (A.8) on the full torus gives
\begin{equation}
 |\log Z_{i,L}'(\beta)+Nf_i^{\mathrm{tr}}(\beta)|
       \le CNe^{-bL}.
 \label{eq:sr-truncated-volume-input}
\end{equation}
The assumptions of Appendix A cover a fixed neighborhood on both sides
of $\beta_c$, even though the main-text Lemma 6.3 is stated on the
ordered side. To identify the stable sides within this same construction, BCT
Lemma A.3 says that the spontaneous magnetization is positive exactly
when $a_o=0$. BCT define their transition point $\beta_0$ as the onset
of positive magnetization, whereas $\beta_c$ above is the crossing
$\psi_o=\psi_d$ of the BKMS free energies. The two coincide: for
$q\ge q_0(d)$ the model has a unique, disordered Gibbs state below the
large-$q$ transition point and $q$ ordered extremal states above it
\cite{KoteckyShlosman1982,LaanaitEtAl1991}, and the BKMS expansion
\eqref{eq:sr-partition-input} is derived from the same Pirogov--Sinai
construction, in which the phase with the smaller $\psi_i$ is the
stable one at each $\beta$. Hence the magnetization vanishes below
$\beta_c$ and is positive above it, so $a_d=0$ below and $a_o=0$
above. Continuity gives $f_o^{\mathrm{tr}}=f_d^{\mathrm{tr}}$ at
coexistence.

The minimum $f^{\mathrm{tr}}=\min_i f_i^{\mathrm{tr}}$ is the physical
thermodynamic random-cluster free energy. This identification can be made
directly in the original convention: BCT Lemma A.1(ii), on the full torus,
bounds each $Z_{i,L}$ above by
$\exp[-Nf^{\mathrm{tr}}+O(Ne^{-bL})]$, while (A.9) supplies the
corresponding lower bound for any minimizing phase. The interface-network
sum (6.26)--(6.27) bounds $Z_{t,L}$ by
$\exp[-Nf^{\mathrm{tr}}-bL^{d-1}]$.
That sum uses only the bounds from Appendix A and the large-$q$ contour
counting inequalities, which hold throughout the fixed neighborhood
under consideration; shrink the neighborhood if necessary so that
$q\le e^{2d\kappa+d}$ there, as required in (6.27).
The sum does not require the ordered phase to minimize.
Combining these estimates with
\eqref{eq:sr-positive-contour-decomposition} identifies the limiting
pressure. Equation~\eqref{eq:sr-rc-normalization} and the stable branch
of~\eqref{eq:sr-partition-input} consequently give
\begin{equation}
 f_i^{\mathrm{tr}}(\beta)=\psi_i(\beta)+d\beta
 \quad\text{on the stable side of phase }i.
 \label{eq:sr-stable-branch-identification}
\end{equation}
We do not identify the two metastable constructions away from their
stable sides or assume differentiability of the minimum cutoff.

\subsection{Weak phase energy limits and nonzero variances}

Applying~\eqref{eq:sr-partition-input} at $\beta_c+t/N$ and using the
boundedness of $U_N/N$ gives
\begin{equation}
 U_N/N\ \Rightarrow\ \frac q{q+1}\delta_{u_o}
                           +\frac1{q+1}\delta_{u_d}.
 \label{eq:sr-macro-limit}
\end{equation}
Indeed the ratios of partition functions are Laplace transforms and
converge for every fixed real $t$ to those of the displayed bounded law.
Define the actual spin events
$B_{o,N}=\{U_N/N<(u_o+u_d)/2\}$ and $B_{d,N}=B_{o,N}^c$.
Their limiting probabilities are the two weights in
\eqref{eq:sr-macro-limit}.

For $s>0$, use~\eqref{eq:sr-partition-input} at
$\beta_c+s/\sqrt N$. The ordered branch dominates, and its quadratic
Taylor expansion gives
\begin{equation}
 \E_{\beta_c}\exp[-s(U_N-Nu_o)/\sqrt N]
       \longrightarrow\frac q{q+1}e^{\sigma_o^2s^2/2}.
 \label{eq:sr-one-sided-laplace}
\end{equation}
The part on $B_{d,N}$ is at most
$e^{-s(u_d-u_o)\sqrt N/2}$ and vanishes. Thus this is also the transform
limit of the positive subprobability measure restricted to $B_{o,N}$.
The disordered phase uses the opposite sign and $\beta_c-s/\sqrt N$.

For clarity, a one-sided transform limit suffices here. Tilt the ordered
subprobability law of $x=(U_N-Nu_o)/\sqrt N$ by $e^{-s_0x}$ for any
fixed $s_0>0$ and normalize it. Ratios of
\eqref{eq:sr-one-sided-laplace} at $s_0-t$ and $s_0$, for $|t|<s_0$,
give convergence of its moment-generating function to that of
$\Normal(-\sigma_o^2s_0,\sigma_o^2)$. The ordinary moment-generating
function convergence theorem gives weak convergence of these tilted
probabilities. Untilting against compactly supported continuous functions
gives vague convergence of the original positive measures. Their known
total masses converge to the recovered Gaussian mass, so vague convergence
is weak convergence. Hence
\begin{equation}
 (U_N-Nu_i)/\sqrt N\mid B_{i,N}
       \Rightarrow\Normal(0,\sigma_i^2),\qquad i=o,d.
 \label{eq:sr-conditional-clt}
\end{equation}
The variance is initially allowed to be zero; nonnegativity follows from
convexity of the limiting log transforms.

Both variances are in fact strictly positive. Here the requirement
$q>2d$ in the choice of $q_0(d)$ is used. Choose an independent
vertex set $S$ of size at least $N/(2d+1)$ and condition on all spins
outside it. The remaining spins are independent. For $x\in S$, let
$n_{x,a}$ count the neighbors of color $a$. Then
$\sum_a n_{x,a}=2d$, and the conditional color probabilities on a
compact inverse-temperature interval with upper endpoint $\beta_{\max}$
obey $p_{x,a}\ge q^{-1}e^{-2d\beta_{\max}}$. If $q>2d$, two neighbor
counts differ by at least one. Since
$\Var_p(n)=\sum_{a<b}p_ap_b(n_a-n_b)^2$, conditional variance gives
\begin{equation}
 \Var_\beta U_N\ge v_*N,\qquad
 v_*:=\frac{e^{-4d\beta_{\max}}}{(2d+1)q^2}>0.
 \label{eq:sr-variance-lower}
\end{equation}
Thus $N^{-1}\log Z_N^{\mathrm{spin}}-v_*\beta^2/2$ is convex.
Its pointwise thermodynamic limit is convex as well. On each open stable
side that limit equals $-\psi_i(\beta)-v_*\beta^2/2$, so
$-\psi_i''\ge v_*$. Continuity of $\psi_i''$ extends the bound to
$\beta_c$, proving $\sigma_i^2\ge v_*>0$. This argument passes through
the stable thermodynamic branches; it does not infer a phase variance
from the much larger total coexistence variance.

\subsection{Uniform moments from positive contour measures}
\label{subsubsec:sr-moments}

We next establish the inward fluctuation control needed for sufficiency.
The finite-volume argument is useful because the minimum in
\eqref{eq:sr-truncation} need not be twice differentiable. The proof
requires only a Lipschitz estimate for its limiting free energy and
then removes every cutoff in a small real temperature interval.

\begin{lemma}[Restricted contour partition derivatives]
\label{lem:sr-restricted-derivatives}
For the fixed sufficiently large $q$ above, there are constants
$\eta,C>0$ such that, for $i=o,d$ and all sufficiently large $L$,
\begin{equation}
 \left|\frac{\dd^2}{\dd\beta^2}\log Z_{i,L}(\beta)\right|\le CN,
 \qquad |\beta-\beta_c|\le\eta/L.
 \label{eq:sr-restricted-second}
\end{equation}
On this interval $Z_{i,L}=Z_{i,L}'$, and
\begin{equation}
 \left.\frac{\dd}{\dd\beta}
 \log[e^{d\beta N}Z_{i,L}(\beta)]\right|_{\beta=\beta_c}
       =-Nu_i+O(\sqrt N).
 \label{eq:sr-restricted-center}
\end{equation}
\end{lemma}
\begin{proof}
We supply the derivative bounds underlying the use of the cluster
expansion. For a general contour interior $\Lambda=\operatorname{Int}\gamma$,
use the positive matching-label sums in BCT (6.13)--(6.14). A summand
indexed by an admissible family $\Gamma$ has logarithm
\begin{equation}
 F_\Gamma(\beta)
 =-e_o(\beta)n_o(\Gamma)-e_d(\beta)n_d(\Gamma)
   -\kappa(\beta)M(\Gamma)+c_\Gamma\log q,
 \label{eq:sr-matching-summand}
\end{equation}
where $n_o+n_d=|\Lambda\cap\Lambda_L|$, the nonnegative integers
$n_o,n_d$ count vertices with each label, and
$M(\Gamma)=\sum_{\gamma'\in\Gamma}m_{\gamma'}$ counts internal contour
intersections. The integer $c_\Gamma$ counts ordered components, with
the constant subtraction of one when the ordered color multiplicity is
removed. All these quantities and the set of admissible families are
independent of $\beta$.

To bound $M(\Gamma)$, recall BCT's construction from cubes of side
$1/2$ and its definition of contour size as the number of
intersections with lattice edges, Section 4 and equation (4.5).
Each contour surface is a union of faces of the half-integer cube
lattice, so it meets a given lattice edge in at most a fixed number
$C_d'$ of points, and pairwise compatible contours are disjoint, so
the intersections counted by $M(\Gamma)$ are distinct points on
distinct edges or on distinct contours. A contour of $\Gamma$ lies in
the closure of $\Lambda$, so the edges it meets either have an
endpoint in $\Lambda\cap\Lambda_L$, of which there are at most
$2d|\Lambda\cap\Lambda_L|$, or cross the bounding contour $\gamma$,
of which there are at most $m_\gamma$ by the definition of contour
size. Consequently
\[
 M(\Gamma)\le C_d'\bigl(2d|\Lambda\cap\Lambda_L|+m_\gamma\bigr)
 \le C_d\bigl(|\Lambda\cap\Lambda_L|+m_\gamma\bigr)
 =:C_d V_\gamma.
\]
The bounded derivatives of $e_o,e_d,\kappa$ now give
$|F_\Gamma'|+|F_\Gamma''|\le C V_\gamma$.
This argument applies to the matching-label sums even when the interior
wraps around the torus; it requires no interpretation as an unrestricted
random-cluster partition function on a simply connected region.

For a positive sum $S=\sum_\Gamma e^{F_\Gamma}$, direct differentiation gives
\[
 (\log S)'=\E_S F_\Gamma',\qquad
 (\log S)''=\E_S F_\Gamma''+\Var_S F_\Gamma'.
\]
Hence $|(\log S)'|\le CV_\gamma$ and
$|(\log S)''|\le C(V_\gamma+V_\gamma^2)$.
Apply these bounds to the numerator and denominator in BCT's activity
ratios (6.18) and (6.22). Since
$V_\gamma=O(m_\gamma^2)$ by~\eqref{eq:sr-contour-geometry}, this
gives the uniform real-parameter bounds
\begin{equation}
 |\partial_\beta\log K_i(\gamma)|\le Cm_\gamma^2,
 \qquad
 |\partial_\beta^2K_i(\gamma)|
                \le Cm_\gamma^4K_i(\gamma).
 \label{eq:sr-activity-derivatives}
\end{equation}
These inequalities do not require the interior phase to be stable.

Here are the precise summability estimates used to differentiate.
For a cluster $\mathcal C$, write
$m(\mathcal C)=\sum_{\gamma\in\mathcal C}m_\gamma$ and let
$\Phi(\mathcal C)$ denote its Ursell coefficient, including the
factor for repeated polymers. The contour counting and incompatibility
estimates (A.5)--(A.6) of~\cite{BorgsChayesTetali2012} give, for some fixed $\rho>0$,
\begin{equation}
 \sum_{\mathcal C}|\Phi(\mathcal C)|
       \prod_{\gamma\in\mathcal C}K_i'(\gamma;\beta)
       e^{\rho m(\mathcal C)}\le CN.
 \label{eq:sr-cluster-majorant}
\end{equation}
The factor $e^{\rho m_\gamma}$ is absorbed by the choice of
$q_0(d)$: the criterion (A.5)--(A.6) holds with a slack
$e^{(c\beta+1)m_\gamma}$ that grows with $q$, so a fixed
$\rho<c\beta+1$ leaves the criterion satisfied on the temperature
interval. To sum over clusters, use the device of BCT's proof of
Lemma A.2 (their Appendix A.3): every contour in $\Lambda_L$ is
incompatible with at least one of the $N$ contours of size two
obtained from configurations with a single disordered edge in a
fixed direction. Every cluster therefore contains a contour
incompatible with one of these $N$ fixed contours, and summing the
incompatibility bound (A.6) over them, with the polynomial factor
absorbed in the exponential slack, proves
\eqref{eq:sr-cluster-majorant}. Overcounting clusters that are
incompatible with several fixed contours only increases the upper
bound.

On any compact real temperature interval, the minimum in
\eqref{eq:sr-truncation} is absolutely continuous. Almost everywhere,
its logarithmic derivative is that of one of its two positive branches;
both obey $|\partial_\beta K_i'|\le Cm_\gamma^2K_i'$.
Termwise first differentiation is therefore justified by
\eqref{eq:sr-cluster-majorant}, since a differentiated cluster is bounded
by $C m(\mathcal C)^2$ times its absolute weight. Equivalently, integrate
these almost-everywhere derivatives over the parameter interval and use
the same majorant to interchange sum and integral. Including the ground
term in~\eqref{eq:sr-polymer-form} gives
$|\partial_\beta\log Z_{i,L}'|\le CN$ almost everywhere. Hence the
finite-volume free energies are uniformly Lipschitz. Their pointwise
limits $f_i^{\mathrm{tr}}$ are Lipschitz with the same constant.

Since $a_i(\beta_c)=0$ and $a_i$ is the difference of two functions
with Lipschitz constant $C$, it follows that
$a_i(\beta)\le 2C|\beta-\beta_c|$. Choose $\eta>0$ small enough that
$2C\eta<c_0\inf\beta$ on the fixed interval. Every torus contour has
diameter at most $L$, so~\eqref{eq:sr-inactive-cutoff-input} now removes
\emph{all} cutoffs on $|\beta-\beta_c|\le\eta/L$.
The finite-volume activities on this interval are their original smooth
positive functions. A second derivative of a cluster product is bounded
by $C m(\mathcal C)^4$ times its absolute weight, by
\eqref{eq:sr-activity-derivatives}. The exponential factor in
\eqref{eq:sr-cluster-majorant} absorbs this polynomial uniformly.
Twice differentiating the uniformly convergent cluster expansion,
and adding the ground term, proves~\eqref{eq:sr-restricted-second}.
This step differentiates no minimum cutoff.

For the center, set
$F_{i,L}(\beta)=\log[e^{d\beta N}Z_{i,L}(\beta)]$ and take
$t_N=\eta/(2\sqrt N)$ with the sign pointing to the stable side of
phase $i$. Since $d\ge2$, $|t_N|\le\eta/L$. Equations
\eqref{eq:sr-truncated-volume-input} and
\eqref{eq:sr-stable-branch-identification} give, both at $\beta_c$ and
at $\beta_c+t_N$,
$F_{i,L}(\beta)=-N\psi_i(\beta)+O(Ne^{-bL})$.
The second-derivative bound and a stable-side finite difference imply
\[
 F_{i,L}'(\beta_c)
 =-N\frac{\psi_i(\beta_c+t_N)-\psi_i(\beta_c)}{t_N}
     +O(N|t_N|)+O(Ne^{-bL}/|t_N|)
 =-Nu_i+O(\sqrt N).
\]
The exponentially small error is negligible for fixed $d$, proving
\eqref{eq:sr-restricted-center}.
\end{proof}

\begin{lemma}[Transfer from bond fluctuations to spin energy]
\label{lem:sr-energy-transfer}
In the Edwards--Sokal joint law at $\beta_c$, the conditional spin
energy on each of the positive bond phase events obeys
\begin{equation}
 \sup_{L\ge L_0,i\in\{o,d\}}
 \E\left[\exp\left(\eta_1\frac{|U_N-Nu_i|}{\sqrt N}\right)
                 \,\middle|\,\mathcal I_{i,L}\right]<\infty
 \label{eq:sr-spin-phase-moment}
\end{equation}
for some fixed $\eta_1>0$ and sufficiently large fixed $L_0$, where
$\mathcal I_{o,L}=\mathcal O_L$ and
$\mathcal I_{d,L}=\mathcal D_L$.
\end{lemma}
\begin{proof}
Introduce the bond fugacity parameter $\lambda=\log(e^\beta-1)$;
then $\dd\beta/\dd\lambda=p$. Put $B=|A|$ and
$M=-U_N$, the number of monochromatic spin bonds. The partition
function $e^{d\beta N}Z_{i,L}(\beta)$ is, up to the constant ordered
multiplicity, a positive sum of $e^{\lambda B}q^{k(A)}$ over the fixed
phase event. Therefore the first and second derivatives of its
\emph{logarithm},
$F_{i,L}(\beta)=\log[e^{d\beta N}Z_{i,L}(\beta)]$, with respect to
$\lambda$ are the conditional mean and variance of $B$.

Lemma~\ref{lem:sr-restricted-derivatives}, the bounded derivatives of
$\beta(\lambda)$, and the elementary $|F_{i,L}'|\le CN$ give a second
$\lambda$ derivative bounded by $CN$ throughout a window of width
$\eta_2/\sqrt N$. At coexistence the first derivative is
$-pNu_i+O(\sqrt N)$, by~\eqref{eq:sr-restricted-center}.
Taylor's theorem for log partition ratios at
$\lambda_c\pm\eta_3/\sqrt N$ therefore proves
\begin{equation}
 \E\left[e^{\eta_3|B+pNu_i|/\sqrt N}
                     \,\middle|\,\mathcal I_{i,L}\right]\le C
 \label{eq:sr-bond-moment}
\end{equation}
for a sufficiently small fixed $\eta_3>0$.

Conditional on the spin configuration, $B$ is Binomial$(M,p)$, since
each monochromatic bond is retained independently with probability $p$.
For a centered Bernoulli variable, the second derivative of its log
moment-generating function is at most $1/4$; integrating twice bounds
that function by $t^2/8$. Conditional independence therefore gives,
for $D=B-pM$,
\begin{equation}
 \E[e^{tD}\mid\sigma]\le e^{dNt^2/8}
 \quad(t\in\R),\qquad
 \E e^{t|D|/\sqrt N}\le2e^{dt^2/8}.
 \label{eq:sr-binomial-noise}
\end{equation}
This estimate remains bounded after conditioning on any event whose
probability is bounded below, at the cost of dividing by that probability.
At coexistence all cutoffs are inactive, and
\eqref{eq:sr-truncated-volume-input} with equal free energies, together
with the tunneling bound, gives
\begin{equation}
 P(\mathcal O_L)\to\frac q{q+1},\qquad
 P(\mathcal D_L)\to\frac1{q+1},\qquad
 P(\mathcal T_L)\le Ce^{-bL^{d-1}}.
 \label{eq:sr-positive-phase-probabilities}
\end{equation}
Thus the conditional noise moments are uniformly bounded. Finally
\[
 |U_N-Nu_i|=|M+Nu_i|
 \le p^{-1}\{|B+pNu_i|+|D|\}.
\]
Cauchy--Schwarz,~\eqref{eq:sr-bond-moment}, and
\eqref{eq:sr-binomial-noise} prove~\eqref{eq:sr-spin-phase-moment}
with $2\eta_1/p\le\eta_3$.
\end{proof}

This transfer uses a conditional binomial fluctuation estimate, not an
identification of occupied bonds with spin energy. In particular,
differentiating the restricted bond partition function gives the bond
count first; the noise estimate is the additional step needed for the
physical Hamiltonian.

\begin{proof}[Proof of Theorem~\ref{thm:sr-bath}]
For necessity, the positive midpoint phase weights,
\eqref{eq:sr-conditional-clt}, and $\sigma_i^2\ge v_*>0$ satisfy
Theorem~\ref{thm:two-scale-necessity}, with gap $N(u_d-u_o)$ and
$s_N=\sqrt N$.

For sufficiency, use the positive decomposition by
$\mathcal O_L,\mathcal D_L,\mathcal T_L$ on the Edwards--Sokal space.
Lemma~\ref{lem:sr-energy-transfer} supplies the physical-energy moment
bound. The exceptional mass is at most $Ce^{-bL^{d-1}}$, while the
largest bath amplification is $\exp[O(N^2/c_N)]$. If
$c_N\gg N^{3/2}$, its logarithm is $o(\sqrt N)=o(L^{d/2})$.
Since $d-1\ge d/2$ for $d\ge2$, the exceptional contribution vanishes.
Theorem~\ref{thm:positive-sufficiency} therefore gives total-variation
convergence on the joint spin--bond space. The likelihood there is
the exact function of $U_N(\sigma)$ in~\eqref{eq:exact-marginal}.
Summing over bonds gives precisely the physical subsystem law, so
projection proves the claimed spin conclusion.
\end{proof}

The argument establishes a small fixed exponential moment on the
$\sqrt N$ scale. It neither asserts a global Gaussian phase tail nor
requires one. Optional independent quadratic momenta preserve the
threshold: their centered Gamma transform adds the moment bound from
Lemma~\ref{lem:mf-exponential}, their mean shifts both centers by
$aN/\beta_c$, and their variance adds $a/\beta_c^2$ to the weak phase
limits. For $a>0$ convolution also yields continuous local Gaussian
energy limits, but no short-range pointwise density envelope is needed
for Theorem~\ref{thm:sr-bath}.
